% Generated draft: requires mathematical and source review.
\documentclass[11pt]{article}
\usepackage[margin=25mm]{geometry}
\usepackage{fontspec}
\setmainfont{Latin Modern Roman}
\usepackage{amsmath,amssymb,mathtools,unicode-math}
\setmathfont{Latin Modern Math}
\usepackage{enumitem,needspace,xurl,hyperref}
\setlength{\parindent}{0pt}
\setlength{\parskip}{5pt}
\setlength{\emergencystretch}{3em}
\newcommand{\N}{\mathbb{N}}
\newcommand{\Z}{\mathbb{Z}}
\newcommand{\Q}{\mathbb{Q}}
\newcommand{\R}{\mathbb{R}}
\newcommand{\C}{\mathbb{C}}
\newcommand{\F}{\mathbb{F}}
\newcommand{\A}{\mathbb{A}}
\newcommand{\PP}{\mathbb{P}}
\newcommand{\E}{\mathbb{E}}
\newcommand{\eps}{\varepsilon}
\newcommand{\dd}{\,\mathrm{d}}
\begin{document}
\section*{UnsolvedMath ID: 30001237}
\textbf{Source problem code:} catalogue entry
\subsection{Definitions and mathematical statement}
The source leaves the ground field and some categorical conventions unspecified. To express its questions, let \(k\) denote a ground field, without imposing a characteristic or algebraic-closure assumption, and let \(X\) be an algebraic variety over \(k\). Write \(D^b(\operatorname{Coh}X)\) for the bounded derived category of coherent sheaves on \(X\), and regard coherent sheaves as complexes concentrated in degree zero.

A coherent sheaf \(E\) is exceptional over \(k\) if
\[
\operatorname{Hom}_{D^b(\operatorname{Coh}X)}(E,E)\cong k,
\qquad
\operatorname{Hom}_{D^b(\operatorname{Coh}X)}(E,E[q])=0
\quad(q\in\Z\setminus\{0\}).
\]
Here the first isomorphism is an isomorphism of \(k\)-algebras. A finite nonempty ordered sequence \((E_1,\ldots,E_m)\) of coherent sheaves is exceptional if every \(E_i\) is exceptional and
\[
\operatorname{Hom}_{D^b(\operatorname{Coh}X)}(E_j,E_i[q])=0
\quad\text{for every }1\leq i<j\leq m\text{ and every }q\in\Z.
\]
It is strongly exceptional if, in addition,
\[
\operatorname{Hom}_{D^b(\operatorname{Coh}X)}(E_i,E_j[q])=0
\quad\text{for every }1\leq i,j\leq m\text{ and every }q\in\Z\setminus\{0\}.
\]
It is full if
\[
\operatorname{thick}(E_1,\ldots,E_m)=D^b(\operatorname{Coh}X),
\]
where \(\operatorname{thick}\) denotes the smallest full triangulated subcategory containing these objects and closed under direct summands. An invertible sheaf is a locally free coherent sheaf of rank one.

\textbf{Source questions.} Does every algebraic variety \(X\) admit a full exceptional sequence of coherent sheaves? Does every such \(X\) admit a strongly exceptional sequence of coherent sheaves? In particular, for every smooth complete rational surface \(X\), do the corresponding sequences exist with every \(E_i\) invertible?

\textbf{Unresolved qualifications.} The source phrase ``a full exceptional sequence, or a strongly exceptional sequence'' does not specify whether the strongly exceptional sequence must also be full. Thus the existence of a strongly exceptional sequence and the existence of a full strongly exceptional sequence must be distinguished; the latter requires both conditions defined above. No required length is supplied for a sequence that is not full. The questions above preserve these alternatives rather than selecting an intended one. The ground field, the precise convention for ``variety,'' and the ambient derived category are not specified by the record; the displayed definitions give an explicit conventional interpretation for review, not verified source conventions. Over a specified field \(k\), \(k\)-rationality of a surface means birationality over \(k\) to \(\PP^2_k\), while geometric rationality means rationality after extension to an algebraic closure. The record does not distinguish these notions. Completeness is retained as stated, and no additional projectivity hypothesis is imposed.
\subsection{Short English statement}
The problem asks about full exceptional sequences and strongly exceptional sequences of coherent sheaves on algebraic varieties, with a special question about sequences of line bundles on smooth complete rational surfaces. It remains unclear whether fullness is required in the strongly exceptional case, and the ground field and rationality convention are unspecified. The supplied label is partially solved; this draft does not independently verify that assessment.
\subsection{Sources}
\begin{itemize}
\item Local UnsolvedMath registry, supplied complete record with permanent ID \(30001237\), title ``Exceptional Sequences on Algebraic Varieties,'' exact source problem code \texttt{catalogue entry}, in the Oberwolfach Reports Open Problems collection. The record gives proposed year \(2009\) and names no proposer.
\item Supplied source citation: \emph{Kommutative Algebra} (2009), Oberwolfach Reports, DOI 10.4171/owr/2009/22, pp.\ 1157--1224. Literature-source label: ``Oberwolfach Report 22, Kommutative Algebra.'' \url{https://doi.org/10.4171/owr/2009/22}. The supplied background describes these reports as documenting mathematical workshops and their open-problem discussions. No problem-specific page or named contributor is supplied.
\item Supplied literature reference: ``Hille-Perling counterexample to King's conjecture.'' \url{https://doi.org/10.1112/S0010437X06002260}.
\item Supplied literature reference: ``Hille-Perling tilting bundles on rational surfaces.'' \url{https://doi.org/10.5802/aif.2860}.
\item Supplied literature assessment, checked on \texttt{2026-08-21}: status ``partially solved.'' The supplied assessment says that existence is established for important classes, including many rational and toric surfaces, that counterexamples exist for full strongly exceptional line-bundle sequences, and that no general existence theorem holds. This is supplied dated curation derived from the OWR record, not independently verified here; its time-sensitive status and cited claims require independent checking.
\item The registry labels the statement \texttt{corrected\_verified} and reports that its mathematical content was checked against the official Oberwolfach report, with notation, context, and question form normalized without changing the source-level claim. This is a supplied verification claim, not a verification performed for this draft.
\end{itemize}
\end{document}
